We build a satellite change-monitoring service for consumers --- a subscriber
pins a neighbourhood and, after each revisit, receives a short written report of
what changed --- from the current AI stack: a promptable segmentation
foundation model (\DetectorName) for open-vocabulary instance detection, CLIP
embeddings for appearance, a spatio-temporal knowledge graph for history, and a
large language model (\LLMName) for the report. Because the deliverable is
prose, we evaluate what it writes, not only what it segments. Two failures that
segmentation metrics cannot see become product failures: a report that invents
construction, and one that miscounts it. Neither is fixed where practitioners
reach first. Adding retrieval or graph aggregation to the language model's
context leaves factuality essentially unchanged; what moves it is the upstream
decision of which detections across two revisits are the same physical object.
A \HeadParamsRound-parameter pairing head, supervised for free from
change-detection masks, reaches \EOnePixelFOne{} pixel F1 on LEVIR-CD against
\AnyChangeFOne{} for the strongest zero-shot peer, and transfers to WHU-CD
unchanged at \ETwoPixelFOne. To measure reports we introduce Change-Fact-Score,
which reduces a report to atomic claims with a deterministic extractor, so no
language model sits inside the metric. It exposes what n-gram overlap hides: a
vision--language model shown both images writes fluent prose while deciding at
chance (\EFourVLMChgAcc\%) whether anything changed. Swapping only the pairing
moves report factuality by \EFiveGain{} points, at \ESixHeadDelta~ms per tile.
